\section{2015级工科数学分析下 B卷}

\subsection{资料来源}
\begin{itemize}
  \item 试题：2015级软件工科数学分析下 B 卷，DOC 正文已抽取；公式对象已按 Word 公式图片逐项恢复。
  \item 公式图片审计：\texttt{tmp/past-exam-equations/2015-B/manifest.csv}。
\end{itemize}

\subsection{试题}
诚信应考，考试作弊将带来严重后果！

华南理工大学本科生期末考试

《工科数学分析》2015--2016 学年第二学期期末考试试卷（B）卷

注意事项：
\begin{enumerate}
  \item 开考前请将密封线内各项信息填写清楚；
  \item 所有答案请直接答在试卷上（或答题纸上）；
  \item 考试形式：闭卷；
  \item 本试卷共 5 个大题，满分 100 分，考试时间 120 分钟。
\end{enumerate}

\subsubsection*{一、填空题（每小题 3 分，共 15 分）}
\begin{enumerate}
  \item 微分方程
  \[
  y''-4y=xe^{2x}
  \]
  的特解形式为 \underline{\hspace{4cm}}。
  \item 设 $z=z(x,y)$ 满足方程
  \[
  x-az=\varphi(y-bz),
  \]
  其中 $\varphi$ 可微，$a,b$ 为常数，则
  \[
  a\frac{\partial z}{\partial x}+b\frac{\partial z}{\partial y}=
  \]
  \underline{\hspace{3cm}}。
  \item 函数
  \[
    f(x,y,z)=\cos(xyz)
  \]
  在点 $P\!\left(\dfrac13,1,\pi\right)$ 处使方向导数取得最大的方向是 \underline{\hspace{4cm}}。
  \item 设 $L:x^2+y^2=a^2$ 取逆时针方向，则
  \[
  \int_L 2xy\,dx+(x^2+8y^2)\,dy=
  \]
  \underline{\hspace{3cm}}。
  \item 设幂级数
  \[
  \sum_{n=0}^{\infty}a_nx^n
  \]
  在 $x=-3$ 条件收敛，则该幂级数的收敛半径为 $R=$ \underline{\hspace{2cm}}。
\end{enumerate}

\subsubsection*{二、计算题（每小题 8 分，共 40 分）}
\begin{enumerate}
  \item 设函数
  \[
  z=f\!\left(xy,\frac{x}{y}\right),
  \]
  其中 $f(\xi,\eta)$ 具有连续的二阶偏导数，求 $\dfrac{\partial^2z}{\partial x\partial y}$。
  \item 计算曲线积分
  \[
  I=\int_\Gamma x^2\,ds,
  \]
  其中 $\Gamma$ 是球面 $x^2+y^2+z^2=a^2$ 与平面 $x+y+z=0$ 的交线。
  \item 设曲线积分
  \[
  \int_L\bigl(f(x)-e^x\bigr)\sin y\,dx-f'(x)\cos y\,dy
  \]
  与路径无关，其中 $f'(x)$ 有一阶的连续导数，且 $f'(0)=0$。

  （1）求 $f'(x)$；（2）计算曲线积分
  \[
  I=\int_{(0,0)}^{(1,1)}\bigl(f(x)-e^x\bigr)\sin y\,dx-f'(x)\cos y\,dy.
  \]

  \item 设 $\Sigma$ 为平面 $y+z=5$ 被柱面 $x^2+y^2=25$ 所截得的部分，计算曲面积分
  \[
  I=\iint_\Sigma (x+y+z)\,dS.
  \]
  \item 设 $\Sigma$ 是上半球面 $z=\sqrt{a^2-x^2-y^2}\ (a>0)$，取上侧，计算曲面积分
  \[
  I=\iint_\Sigma (x^3-xy^2)\,dy\,dz+(y^3-yz^2)\,dz\,dx+(z^3-zx^2)\,dx\,dy.
  \]
\end{enumerate}

\subsubsection*{三、解答题（每小题 9 分，共 18 分）}
\begin{enumerate}
  \item 将函数
  \[
  \frac{d}{dx}\left(\frac{e^x-1}{x}\right)
  \]
  展开成 $x$ 的幂级数，并求数项级数
  \[
  \sum_{n=1}^{\infty}\frac{n}{(n+1)!}
  \]
  的和。
  \item 求微分方程
  \[
  y'\tan x=y\ln y
  \]
  的通解。
\end{enumerate}

\subsubsection*{四、证明题（每小题 9 分，共 18 分）}
\begin{enumerate}
  \item 设 $f(x)\in C[a,b]$，证明：
  \[
  \int_a^b dx\int_a^x f(y)\,dy=\int_a^b f(y)(b-y)\,dy.
  \]
  \item 证明：函数项级数
  \[
  \sum_{n=1}^{\infty}\frac{(-1)^n(1-e^{-nx})}{n^2+x^2}
  \]
  在 $[0,+\infty)$ 上一致收敛。
\end{enumerate}

\subsubsection*{五、应用题（本题 9 分）}
从斜边之长为 $a$ 的一切直角三角形中，求有最大周长的直角三角形。

\subsection{答案与解析}

\subsubsection*{一、填空题}
\begin{enumerate}
  \item 特征方程 $r^2-4=0$ 有根 $r=\pm2$，右端为一次多项式乘 $e^{2x}$，且 $2$ 是一重特征根，故特解形式为
  \[
  y^*=xe^{2x}(Ax+B).
  \]
  \item 记 $u=y-bz$。对 $x,y$ 分别求偏导：
  \[
  1-az_x=-b\varphi'(u)z_x,\qquad -az_y=\varphi'(u)(1-bz_y).
  \]
  化简得
  \[
  a z_x+b z_y=1.
  \]
  \item 最大方向为梯度方向。因
  \[
  \nabla f=-(yz,xz,xy)\sin(xyz),
  \]
  在 $P\!\left(\dfrac13,1,\pi\right)$ 处有
  \[
  \nabla f(P)=-\frac{\sqrt3}{6}(3\pi,\pi,1).
  \]
  故所求单位方向可取
  \[
  -\frac{(3\pi,\pi,1)}{\sqrt{10\pi^2+1}}.
  \]
  \item 由 Green 公式，
  \[
  \int_L 2xy\,dx+(x^2+8y^2)\,dy
  =\iint_D(2x-2x)\,dA=0.
  \]
  \item 幂级数在 $x=-3$ 条件收敛，说明 $x=-3$ 是收敛区间端点，故收敛半径为
  \[
  R=3.
  \]
\end{enumerate}

\subsubsection*{二、计算题}
\begin{enumerate}
  \item 令 $\xi=xy,\ \eta=\dfrac{x}{y}$。先求
  \[
  z_x=yf_\xi+\frac1y f_\eta.
  \]
  再对 $y$ 求偏导，混合项相消，得
  \[
  z_{xy}=f_\xi+xyf_{\xi\xi}-\frac{x}{y^3}f_{\eta\eta}-\frac1{y^2}f_\eta,
  \]
  其中各偏导均在 $\left(xy,\dfrac{x}{y}\right)$ 处取值。
  \item $\Gamma$ 是半径为 $a$ 的大圆，所在平面的单位法向量为
  \[
  \mathbf n=\frac{(1,1,1)}{\sqrt3}.
  \]
  单位向量 $\mathbf e_x$ 在该平面内的投影长度平方为 $1-\dfrac13=\dfrac23$。故圆周上 $x^2$ 的平均值为
  \[
  \frac{a^2}{2}\cdot\frac23=\frac{a^2}{3}.
  \]
  因此
  \[
  I=2\pi a\cdot\frac{a^2}{3}=\frac{2\pi a^3}{3}.
  \]
  \item 记
  \[
  P=(f(x)-e^x)\sin y,\qquad Q=-f'(x)\cos y.
  \]
  路径无关给出 $P_y=Q_x$，即
  \[
  f''(x)+f(x)=e^x.
  \]
  通解为
  \[
  f(x)=C_1\cos x+C_2\sin x+\frac12e^x.
  \]
  又 $f'(0)=0$，得 $C_2=-\dfrac12$，于是
  \[
  f'(x)=-C_1\sin x-\frac12\cos x+\frac12e^x.
  \]
  因为
  \[
  (f-e^x)\sin y\,dx-f'(x)\cos y\,dy=d\bigl[-f'(x)\sin y\bigr],
  \]
  所以
  \[
  I=-f'(1)\sin1
  =C_1\sin^2 1+\frac12(\cos1-e)\sin1.
  \]
  原题只给出 $f'(0)=0$，只能确定 $C_2=-\dfrac12$，不能唯一确定 $C_1$。因此在现有题面条件下，积分值应保留为上述含任意常数 $C_1$ 的表达式；若要得到唯一数值，还需另给一个关于 $f$ 或 $f'$ 的独立条件。
  \item 将平面写为 $z=5-y$，则
  \[
  dS=\sqrt2\,dx\,dy,\qquad x+y+z=x+5.
  \]
  投影区域为圆盘 $D:x^2+y^2\le25$，故
  \[
  I=\sqrt2\iint_D(x+5)\,dx\,dy
  =125\sqrt2\,\pi.
  \]
  \item 令
  \[
  \mathbf F=(x^3-xy^2,\ y^3-yz^2,\ z^3-zx^2).
  \]
  则
  \[
  \nabla\cdot\mathbf F
  =(3x^2-y^2)+(3y^2-z^2)+(3z^2-x^2)
  =2(x^2+y^2+z^2).
  \]
  上半球面与底面圆盘围成上半球体，底面上 $z=0$，第三分量通量为 $0$。于是
  \[
  I=\iiint_\Omega 2(x^2+y^2+z^2)\,dV
  =2\cdot 2\pi\int_0^a r^4\,dr
  =\frac{4\pi a^5}{5}.
  \]
\end{enumerate}

\subsubsection*{三、解答题}
\begin{enumerate}
  \item 先写
  \[
  \frac{e^x-1}{x}=\sum_{n=0}^{\infty}\frac{x^n}{(n+1)!}.
  \]
  逐项求导得
  \[
  \frac{d}{dx}\left(\frac{e^x-1}{x}\right)
  =\sum_{n=1}^{\infty}\frac{n}{(n+1)!}x^{n-1}.
  \]
  令 $x=1$，也可由闭式
  \[
  \frac{d}{dx}\left(\frac{e^x-1}{x}\right)
  =\frac{(x-1)e^x+1}{x^2}
  \]
  得
  \[
  \sum_{n=1}^{\infty}\frac{n}{(n+1)!}=1.
  \]
  \item 设 $y>0$，方程化为
  \[
  \frac{dy}{y\ln y}=\cot x\,dx.
  \]
  积分得
  \[
  \ln|\ln y|=\ln|\sin x|+C,
  \]
  即
  \[
  \ln y=C\sin x,\qquad y=e^{C\sin x}.
  \]
  其中 $C=0$ 时包含常值解 $y=1$。
\end{enumerate}

\subsubsection*{四、证明题}
\begin{enumerate}
  \item 左端积分区域为 $a\le y\le x\le b$。交换积分次序得
  \[
  \int_a^b dx\int_a^x f(y)\,dy
  =\int_a^b dy\int_y^b f(y)\,dx
  =\int_a^b f(y)(b-y)\,dy.
  \]
  \item 对 $x\in[0,+\infty)$，有
  \[
  0\le 1-e^{-nx}\le1,\qquad n^2+x^2\ge n^2.
  \]
  因而
  \[
  \left|\frac{(-1)^n(1-e^{-nx})}{n^2+x^2}\right|\le\frac1{n^2}.
  \]
  由 $\sum_{n=1}^{\infty}\dfrac1{n^2}$ 收敛及 Weierstrass 判别法，原级数在 $[0,+\infty)$ 上一致收敛。
\end{enumerate}

\subsubsection*{五、应用题}
设直角三角形两直角边为 $x,y>0$，则
\[
x^2+y^2=a^2,\qquad P=a+x+y.
\]
在圆弧约束下，$x+y$ 最大时 $x=y$，故
\[
x=y=\frac{a}{\sqrt2}.
\]
因此所求三角形为等腰直角三角形，最大周长为
\[
P_{\max}=a+\sqrt2\,a=(1+\sqrt2)a.
\]
